\section{试验台规格：如何复现}
\label{sec:dishspec}

下面汇总了重新搭建这样一个试验台所需的全部内容：设备、回路、输入编码、输出读取、反馈和实验方案。每个参数都注明了来源。“论文”表示数值取自文献~\cite{Kagan2022} 的正文。“自选”表示论文中没有给出，复现时需要自行设定；我们给出默认值，并说明如何检验。

\subsection*{设备与培养物}

\begin{center}
\footnotesize
\setlength{\tabcolsep}{4pt}
\renewcommand{\arraystretch}{1.15}
\begin{tabular}{>{\raggedright}p{4.0cm}>{\raggedright}p{6.9cm}>{\raggedright\arraybackslash}p{1.6cm}}
\toprule
参数 & 数值 & 来源 \\
\midrule
芯片 & MaxOne 阵列：$8$~mm$^2$ 区域上 $26\,000$ 个铂电极 & 论文 \\
记录 & 最多同时 $1024$ 个通道，每秒 $20\,000$ 个采样 & 论文 \\
刺激 & 技术上最多 $32$ 个电极，实验中用 $8$ 个独立电极 & 论文 \\
细胞 & 小鼠胚胎皮层；由干细胞获得的人类神经元（两种培养方法）；HEK293T 细胞（非神经元）作对照 & 论文 \\
接种 & 每块芯片约 $10^6$ 个细胞 & 论文 \\
实验时的培养龄 & 小鼠：从 $\sim10$ 天起；人类：视培养方法为 $14$--$24$ 天或 $\sim80$ 天 & 论文 \\
\bottomrule
\end{tabular}
\end{center}

\subsection*{回路与时间}

回路以 $10\ms$（$200$ 个采样）为窗口运行：每个窗口内统计运动区电极上的脉冲，更新游戏，并发出下一次刺激。从培养物中的脉冲到响应刺激的延迟约 $5\ms$。游戏期间，每个电极的信号经过截止频率 $100$~Hz 的二阶贝塞尔高通滤波器；信号的绝对值再经 $1$~Hz 的一阶贝塞尔低通滤波器平滑，脉冲阈值与这个平滑后的绝对值成正比。论文中六倍背景标准差的阈值用于单独的活动测量，而不是游戏。为了不把刺激误算为培养物的响应，当同时出现超过 $15$ 个大于 $75$~mV 的大脉冲时，所有信号都被屏蔽。以上均来自论文。

\subsection*{球的编码}

\begin{center}
\footnotesize
\setlength{\tabcolsep}{4pt}
\renewcommand{\arraystretch}{1.15}
\begin{tabular}{>{\raggedright}p{3.4cm}>{\raggedright}p{7.5cm}>{\raggedright\arraybackslash}p{1.6cm}}
\toprule
参数 & 数值 & 来源 \\
\midrule
感觉区 & $8$ 个电极，排列使相邻电极表示球的相邻位置 & 论文 \\
位置编码 & 电极编号 $k=1,\dots,8$ 表示球相对于球拍中心的位置 & 论文 \\
哪个坐标 & 球的竖直位置；复现时取相对于球拍的位置，$y-p\in[-\tfrac12,\tfrac12]$ 时 $k=\lceil 8(y-p+\tfrac12)\rceil$ & 论文；公式为自选 \\
频率编码 & $4$ 到 $40$~Hz 的刺激频率表示到球拍的距离 & 论文 \\
标度方向 & 球在对面墙边时 $4$~Hz，线性增加到球在球拍一侧墙边时的 $40$~Hz：$f=4+36\,(1-x/L)$，$x$ 为到球拍一侧墙的距离，$L$ 为场地宽度 & 论文 \\
脉冲波形 & 短的矩形双相脉冲：先正相，后负相 & 论文 \\
幅度 & $75$~mV；选定这一幅度是为了不强迫受抑制的细胞发放 & 论文 \\
相位时长 & 未给出；采用刺激系统的标准设置，实验期间不改变 & 自选 \\
\bottomrule
\end{tabular}
\end{center}

\subsection*{读取球拍}

论文中有两个运动区：第一个区的脉冲让球拍上移，第二个区的让它下移；两个区的贡献经过加权，以补偿细胞密度和活动水平的差异~\cite{Kagan2022}。确切公式没有给出。复现时我们采用规则~\eqref{eq:dishreadout}，每个 $10\ms$ 窗口 $\Delta p=c\,(n_\uparrow-n_\downarrow)$，用权重按静息平均活动把两个区拉平（自选）：$n_\uparrow$ 和 $n_\downarrow$ 各自除以游戏前测得的本区每窗口平均计数。系数 $c$ 取为使一个平均计数的差每窗口把球拍移动场地高度的 $1\%$（自选）；这样，一个区持续占优时，球拍 $1$~s 就能走完整个场地。

论文正文没有用坐标给出各区在芯片上的位置，只有一张示意图。复现时只需三组互不重叠的电极：中间八个感觉电极，两侧各一组记录电极，一组为“上”，一组为“下”（自选）。

\subsection*{游戏}

论文没有给出场地尺寸、球拍长度和球速；只说球以恒定速度飞行并在墙上反弹。复现时（自选）：场地 $1\times1$，球拍在右侧墙上，长 $0.2$（$|y-p|<0.1$ 即击中，与模型~\texttt{models/dishbrain\_model.py} 一致），球 $1$~s 飞越场地。一个回合中一次都没接住就未接住，记为“ace”；连续击球超过三次的回合记为“长回合”（论文）。

\subsection*{反馈}

\begin{center}
\footnotesize
\setlength{\tabcolsep}{4pt}
\renewcommand{\arraystretch}{1.15}
\begin{tabular}{>{\raggedright}p{2.6cm}>{\raggedright}p{8.3cm}>{\raggedright\arraybackslash}p{1.6cm}}
\toprule
事件 & 施加的刺激 & 来源 \\
\midrule
击中 & 可预测刺激：全部 $8$ 个感觉电极同时，$75$~mV，$100$~Hz，$100\ms$；其间暂停通常的球位置刺激 & 论文 \\
未击中 & 不可预测刺激：$150$~mV，$5$~Hz，$4$~s，在随机时刻施加到从 $8$ 个电极中随机选出的电极上；随后 $4$~s 无刺激，球沿随机方向重新发出 & 论文 \\
随机规律 & 未给出；复现时把脉冲时刻取为平均频率 $5$~Hz 的泊松流 & 自选 \\
\bottomrule
\end{tabular}
\end{center}

\subsection*{实验方案与对照}

一次会话持续 $20$~min；比较前 $5$~min 和后 $15$~min（论文）。三项结果指标：平均回合长度、ace 比例和长回合比例。实验条件（论文）：
\begin{itemize}
\item \emph{刺激}：完整回路，如上所述；
\item \emph{安静}：以同样长的无任何刺激时段代替击中或未击中刺激，然后球沿随机方向发出；
\item \emph{无反馈}：未接住后球直接反弹继续飞行，球位置刺激继续进行；
\item \emph{静息}：游戏进行、球拍随活动移动，但完全没有刺激；
\item \emph{对照}：只有培养基的芯片；HEK293T 细胞；“计算机里的培养物”，即用仿真代替细胞。
\end{itemize}
主系列的样本量：小鼠培养物 $9$ 个（$101$ 次会话），人类培养物 $11$ 个（$138$），只有培养基的芯片 $6$ 块（$80$），静息培养物 $20$ 个（$42$），仿真 $3$ 个（$38$），共 $399$ 次会话。反馈系列：$3$ 天、每天 $3$ 次会话，共 $486$ 次会话，条件为“刺激”“安静”和“无反馈”（$n=27$、$17$ 和 $15$）。平均回合长度从前 $5$ 分钟到后 $15$ 分钟的增长只出现在活的培养物中。在反馈系列中，随时间的显著增长只出现在“刺激”条件下；会话结束时，“安静”条件仍优于“无反馈”，但效应较小；而在三天里跨会话的稳定学习并未出现~\cite{Kagan2022}。

\subsection*{试验台循环的步骤}

每隔 $10\ms$，试验台计算机执行以下操作。
\begin{enumerate}
\item 读取上一窗口内两个运动区的脉冲数 $n_\uparrow$、$n_\downarrow$，并用各区静息平均计数归一化。
\item 把球拍移动 $\Delta p=c\,(n_\uparrow-n_\downarrow)$，并限制在场地边界内。
\item 把球移动一步；在上、下、左三面墙上反弹。
\item 如果球到达右墙：若 $|y-p|<0.1$ 则为击中，回合计数加一，施加击中刺激（$100\ms$）；否则为未击中，记录该回合，施加未击中刺激（$4$~s），再停顿 $4$~s，球重新发出。
\item 否则，按球的竖直位置选择电极 $k$，按到球拍一侧墙的距离选择频率 $f$，在下一个窗口之前向电极 $k$ 施加频率为 $f$、幅度 $75$~mV 的双相脉冲。
\end{enumerate}
这些足以搭建一个与已发表试验台兼容的装置，并检验本章的核心问题：教会培养物的是可预测性，还是仅仅是对击中与未击中响应之间的差别。为此，应在论文的三个条件之外加上论文中没有的第四个：未接住后施加与不可预测刺激时长和能量相同的\emph{可预测}刺激。如果培养物这样也能学会，关键就在于差别，而不在于可预测性。
